% Created 2018-10-18 Thu 15:27
% Intended LaTeX compiler: pdflatex
\documentclass[11pt]{article}
\usepackage[utf8]{inputenc}
\usepackage[T1]{fontenc}
\usepackage{graphicx}
\usepackage{grffile}
\usepackage{longtable}
\usepackage{wrapfig}
\usepackage{rotating}
\usepackage[normalem]{ulem}
\usepackage{amsmath}
\usepackage{textcomp}
\usepackage{amssymb}
\usepackage{capt-of}
\usepackage{hyperref}
\pagenumbering{roman}
\usepackage[margin=0.80in]{geometry}
\usepackage[backend=biber, style=ieee, url=false, dashed=false]{biblatex}
\usepackage{float}
\usepackage[super,negative]{nth}
\usepackage[capitalise]{cleveref}
\usepackage{pst-node,transparent,ragged2e}
\addbibresource{/home/nlisa/.spacemacs.d/org-files/bibliography.bib}
\DeclareFieldFormat[inproceedings]{citetitle}{\textit{#1}}
\DeclareFieldFormat[inproceedings]{title}{\textit{#1}}
\DeclareFieldFormat[misc]{citetitle}{#1}
\DeclareFieldFormat[misc]{title}{#1}
\renewcommand*{\bibpagespunct}{%
\ifentrytype{inproceedings}
{\addspace}
{\addcomma\space}}
\AtEveryCitekey{\ifuseauthor{}{\clearname{author}}}
\AtEveryBibitem{\ifuseauthor{}{\clearname{author}}}
\SpecialCoor
\author{Nyameko Lisa}
\date{19 October 2018}
\title{IPSP078 - Legal Aspects of Electronic Commerce\\\medskip
\large Portfolio Exam: 691859}
\hypersetup{
 pdfauthor={Nyameko Lisa},
 pdftitle={IPSP078 - Legal Aspects of Electronic Commerce},
 pdfkeywords={},
 pdfsubject={},
 pdfcreator={Emacs 26.1 (Org mode 9.1.14)},
 pdflang={English}}
\begin{document}

\maketitle
\addvspace{110pt}
\centering{
\pnode(0.5\textwidth,-0.5\textheight){thisCenter}
\rput(thisCenter){%\transparent{0.25}
\includegraphics[width=2.7in]{/home/nlisa/course/llb/wipo-unisa/UNISACoatofArms.eps}}}

\justifying
\addvspace{110pt}
\section*{Declaration}
\label{sec:orgb4520eb}
I know that plagiarism is to use someone else’s work and pass it off as my own.\\
I know that plagiarism is wrong.\\
I confirm that this assignment is my own work.\\
I have acknowledged in the bibliography accompanying the assignment all the sources that I have used.\\
I have not directly copied without acknowledgment anything from the Internet or from any other source.\\
I have indicated every quotation and citation in a footnote or bracket linked to that quotation.\\
I have not allowed anyone else to copy my work and to pass it off as their own work.\\
I understand that if any unacknowledged copying whatsoever appears in my assignment I will receive zero per cent for the assignment.\\
I am aware of the UNISA policy on plagiarism and understand that disciplinary proceedings can be instituted against me by UNISA if I contravene this policy.\\
I indicate my understanding and acceptance of this declaration by
entering my name here under:
\begin{itemize}
\item Name: \textbf{Nyameko Lisa} (Student Number: \textbf{7874-909-3})
\end{itemize}

\subsection*{NOTE}
\label{sec:orgcc148e9}
Please note that footnotes will be denoted as \footnote{This is a footnote.} and will
appear at the bottom of the page.\\
References will be denoted by \cite{rsa02_elect_comm_trans_act} and a numerically
organized citation list will appear at the end of the document.
\newpage

\pagenumbering{arabic}
\section{Ofentse approaches you and wants your advise on the following: [25]}
\label{sec:orgdea6a81}

As per the provisions of [section 1]\cite{rsa02_elect_comm_trans_act} together
with [section 1(q)]\cite{rsa12_elect_comm_trans_amend_bill} the data base will be
understood to constitute either \textbf{`data'} or a \textbf{`data message'}, Ofentse as
the compiler of the database will be understood to be the \textbf{`data controller'}
and his clients whose personal information constitute the content of the
database will be understood to be \textbf{`data subjects'}.

\subsection{The requirements that the databases must meet in order to enjoy copyright protection. (10)}
\label{sec:orgb7d230c}

As per the provisions of [section 1(1)]\cite{rsa78_copyrightact} in the context of
copyright law a database i.e. \emph{tables and compilations of data stored or
embodied in a computer or medium used in conjunction with a computer,} will be
considered to be a \textbf{`literary work'} and as a such it will be eligible for
copyright per the provisions of [section 2(1)(a)]\cite{rsa78_copyrightact}
provided Ofentse can demonstrate that his database is indeed original, and
finally as per [section 2(2)]\cite{rsa78_copyrightact} that the database can be
recorded in digital data. Furthermore copyright shall be conferred to the
database if Ofentse can demonstrate that he is a \textbf{`qualified person'}, i.e. that
he is a South African citizen or is domiciled or resident in the Republic
[section 3(1)(a)]\cite{rsa78_copyrightact}. \\

As a Paris Convention Member Statue the provisions of [article
9(1)]\cite{wto17_trips}, [section 5]\cite{rsa78_copyrightact} [article
1(2)]\cite{eurlex00_elec_commerce}, [article 5(2)]\cite{wipo86_berne} and [article
1(4)]\cite{wipo96_copyright_treaty} dictate that the statues of the Republic shall
govern international copyright protection as it pertains to the Internet.

\subsection{The provisions of the Electronic Communications and Transactions Act 2002 which limits the liability of the Internet service providers. (15)}
\label{sec:org140a1a1}

As per the provisions of the ETC Act \cite{rsa02_elect_comm_trans_act}, a service
provider is not liable for damages arising from said service provider for
providing access to information system services or wireless applications
services [section 70]\cite{rsa02_elect_comm_trans_act} together with [section
36]\cite{rsa12_elect_comm_trans_amend_bill}, or for operating facilities as a
\emph{`mere conduit'} [section 73]\cite{rsa02_elect_comm_trans_act} together with
[section 38]\cite{rsa12_elect_comm_trans_amend_bill}, \emph{`caching purposes'}
[section 74]\cite{rsa02_elect_comm_trans_act} together with [section
39]\cite{rsa12_elect_comm_trans_amend_bill} and \emph{`hosting services'} [section
75]\cite{rsa02_elect_comm_trans_act}. The eligibility for the limitation of
liability of a service provider is dependent on their membership with a
recognized Representative body as per [sections 71 and
72]\cite{rsa02_elect_comm_trans_act} together with [section
37]\cite{rsa12_elect_comm_trans_amend_bill} Moreover the Act explicitly excludes
service providers' obligation from monitoring data, as per [section
78]\cite{rsa02_elect_comm_trans_act}.

\section{Advise the Committee on the following issues [25]:}
\label{sec:org0775555}

\subsection{Does EduCate's marketing practice constitute spam according to South African law? If so, discuss those legislative provisions and state whether they adequately protect consumers? [15]}
\label{sec:org4449c19}

Given that the recipients email addresses were retrieved from public websites,
and that there have been consumers who've lodged complaints against EduCate CC,
EduCate's bulk email newsletters constitute commercially motivated
direct-marketing by means of unsolicited electronic communications, thus
amounting to spam
\cite{hermann14_direct_marketing_vs_spam,tladi08_reg_unsol_comm}. As per the
provisions of [section 45]\cite{rsa02_elect_comm_trans_act} together with [section
23]\cite{rsa12_elect_comm_trans_amend_bill} recipients of unsolicited
communications are able to opt-out of future communications and may request
information on where their personal information and contact details were
obtained. Moreover as per the provisions of [section
45(3,4)]\cite{rsa02_elect_comm_trans_act} together with [section
23]\cite{rsa12_elect_comm_trans_amend_bill}, anyone failing to comply or
persisting to send unsolicited commercial communications to a person whose
advised that such communication is unwelcome, is guilty of an offense and liable
on conviction to penalties.\\

As per the provisions of [section 11(1, 2)]\cite{rsa08_cpa} every
person has the right to privacy which includes the right to restrict unwanted
direct marketing through retrospective refusal to accept or request
discontinuance or preemptive blockage of said unsolicited messages. Moreover as
per the provisions of [section 11(3 - 5)]\cite{rsa08_cpa}
consumers may register either a general or specific preemptive block against any
communications primarily for the purpose of direct marketing.\\

As per the provisions of [section 69(1)]\cite{rsa13_popi} stipulates that for the
purposes of direct marketing by any form of electronic communication, the
processing of personal information of a data subject is prohibited unless the
data subject has given their consent. Moreover as per [subsections (2, 3) of
section 69]\cite{rsa13_popi} a responsible party may approach existing customers
or prospective consumers on a single occasion provided that said consumers have
not previously withheld such consent. Lastly as per the provisions of [section
69(4)]\cite{rsa13_popi} stipulates that a responsible party must in their
communications include identifying details of either the sender or party on
behalf of whom the information has been sent, and a mechanism for the cessation
of further messages.

\subsection{Choose one country that regulates spam in its jurisdiction and give a detailed discussion on the country's initiatives. Also note what SA can learn from that jurisdiction regarding the issues of spam. (10)}
\label{sec:org9717a87}

Canada's Anti-Spam Legislation \cite{ca10_anti_spam_act} was enacted on 15
December 2010 and its purpose is to promote the efficiency and adaptability of
the Canadian economy [section 3]\cite{ca10_anti_spam_act} by regulating a range of
online activities including spam messaging, spyware distribution, fishing,
fraudulent as well as other deceptive and harmful online threats and practices.
The legislation adopts a holistic approach in encompassing through
implementation within this Act amendments to their corresponding legislation:
the Canadian Radio-Television and Telecommunications Commission Act
\cite{ca85_radio_tv_telecomms_act}, the Canadian Competition Act
\cite{ca85_comp_act}, the Canadian Personal Information Protection and Electronic
Documents Act \cite{ca00_popi_act} and the Canadian Telecommunications Act
\cite{ca93_telecomms_act}.\\

The introductory [sections 1 - 5]\cite{ca10_anti_spam_act} deal with the
interpretive definitions; conflicting provisions and the precedence of this Act;
the purpose of the Act, scope and applicability of the Act respectively.
[Sections 6 - 13]\cite{ca10_anti_spam_act} detail the \textbf{Requirements and
Prohibitions} related to electronic messages during the conduct of commercial
activities, namely: unsolicited electronic messages [section
6]\cite{ca10_anti_spam_act}, alteration of transmission data [section
7]\cite{ca10_anti_spam_act}, installation of computer programs [section
8]\cite{ca10_anti_spam_act}, the jurisdiction where the contravention of the
corresponding sections applies [sections 9, 12]\cite{ca10_anti_spam_act}, the
requirements for express consent to perform regulated and prohibited actions
[section 10]\cite{ca10_anti_spam_act}, mandatory provision of an unsubscribe
mechanism with respect to unsolicited messages [section
11]\cite{ca10_anti_spam_act} and lastly how the burden of proof vests with the
alleged perpetrator of a prohibited activity to demonstrate that they indeed had
express consent to perform said activity [section
13]\cite{ca10_anti_spam_act}.\\

\textbf{Administrative Monetary Penalties} are articulated in [sections
14 - 40]\cite{ca10_anti_spam_act}, namely: the appointment of a designated authority
[section 14]\cite{ca10_anti_spam_act}; their authority to demand preservation of
and notice to produce infringing electronic data [section
15 - 18]\cite{ca10_anti_spam_act}; issuance and execution of warrants, including
details regrading use of force and entry onto private property [section
19]\cite{ca10_anti_spam_act}; notification and determination of responsibility for
the contravention of the prohibitions amounting to violations for which one may
be found liable to an administrative monetary penalty to the value of \$1,000,000
in the case of an individual and \$10,000,000 in the case of any other person,
are considered in [sections 20 - 26]\cite{ca10_anti_spam_act}, these sections also
detail how the penalties are not intended to be punitive, but rather they are
intended to encourage compliance; the appeals process to the Federal Court of
Appeals, the recovery of penalties and other amounts, and the rules regarding
violations are the purview of [sections 27 - 33]\cite{ca10_anti_spam_act}, including
\emph{vicarious liability} where an employer may be liable for violations committed
by their employee during the course of their employment; general provisions
regarding judicial powers, rules of procedure, evidence, publication and
enforcement are dealt with in [sections 34 - 40]\cite{ca10_anti_spam_act}.\\

Notice and issuance of an \textbf{Injunction} is outlined in [section
41]\cite{ca10_anti_spam_act}, whilst matters relating to \textbf{Offenses} are
articulated in [section 42 - 46]\cite{ca10_anti_spam_act}. These include offenses
resulting from non-compliance [section 42]\cite{ca10_anti_spam_act} with respect
to refusal or failure to comply with either a preservation demand [section
15]\cite{ca10_anti_spam_act}, presentation notification [section
17]\cite{ca10_anti_spam_act}, and/or a warrant issued as per [section
19(4)]\cite{ca10_anti_spam_act}; obstruction and falsification of information
[section 43]; implicit and vicarious liability of directors and officers of
corporations and employers in general [sections 44 \& 45]\cite{ca10_anti_spam_act};
with the monetary details for offenses detailed in [section
46]\cite{ca10_anti_spam_act}.\\

With regards to actual litigation and the \textbf{Private Right of Action} [sections
47 - 49]\cite{ca10_anti_spam_act} deal with the application process; [sections
50 - 51]\cite{ca10_anti_spam_act} provide details for the court hearing and their
respective orders; and [sections 52 - 55]\cite{ca10_anti_spam_act} describe the
rules about contraventions and re-viewable conduct. The remaining sections
details statutory regulations, as well as parliamentary and ministerial roles
and responsibilities. Where [sections 56 - 61]\cite{ca10_anti_spam_act} described
the process for the \textbf{Consultation and Disclosure of Information}, be it by an
organization [section 56 - 57]\cite{ca10_anti_spam_act}; by a Commission [sections
58 - 59]\cite{ca10_anti_spam_act}; by the government of a foreign state [section
60]\cite{ca10_anti_spam_act}; and the reports to the Ministry of Industry [section
61]\cite{ca10_anti_spam_act}.\\

The Canadian Anti-Spam Law \cite{ca10_anti_spam_act} is arguably one of the most
stringent in the world and should conceivably have a significant impact on
reducing spam, in that it constitutes specific legislation with a targeted
purpose [section 3]\cite{ca10_anti_spam_act}, as opposed to a fragmented,
piecemeal collection of a limited amount of provisions
\cite{hermann14_direct_marketing_vs_spam,tladi08_reg_unsol_comm} which are secondary to the aims of the
statues within which they are contained
\cite{rsa02_elect_comm_trans_act,rsa12_elect_comm_trans_amend_bill,rsa08_cpa,rsa13_popi}
as is the case within the Republic.

\section{Advise Kate on the following: [25]}
\label{sec:org910ac7e}

\subsection{Whether the requirements of an electronic signature have been met by clicking on the icon ``submit order''. (10)}
\label{sec:org1527dfb}

\subsection{Does South African law recognize this type of transaction, and what are the legal implications thereof? (10)}
\label{sec:org2293a88}

\subsection{Can she cancel the contract? (5)}
\label{sec:org222f69f}

\section{Discuss the following questions regarding the term permanent establishment (PE): [25]}
\label{sec:org0497c7e}

\subsection{Describe what the term PE means. (5)}
\label{sec:orgdd29933}

A \textbf{permanent establishment (PE)} is a fixed place of business through which the
business of an enterprise is wholly or partially carried on, [article
5(1)]\cite{oecd17_model_tax}, generally giving rise to income or value added tax
(VAT) liability [article 7(1)]\cite{oecd17_model_tax} in a particular jurisdiction
resulting from the bilateral tax treaties negotiated for the treatment of
cross-border commerce. A (PE) is created in terms of a \textbf{Double Taxation
Agreement (DTA)}, falling under the ambit of legislation to be administered by
an authorized person or commissioner as described in [schedule
1]\cite{rsa97_sars_act}, and details the requirements a non-resident enterprise
must meet under which profits generated through commercial actives in a foreign
state may subsequently be taxed in said foreign state. As per the provisions of
[section 1]\cite{rsa62_income_tax}, which directly refers to [article
5(2)]\cite{oecd17_model_tax}, examples of a (PE) include: a place of management; a
branch; an office; a factory; a workshop; a mine, an oil or gas well, a quarry
or any other place of extraction of natural resources; while a number of
exclusions are detailed in [article 5(4)]\cite{oecd17_model_tax}. A (PE) must satisfy the following prerequisites:
\begin{itemize}
\item \textbf{place-of-business test} - there must necessarily exist a distinct premises,
equipment or machinery,
\item \textbf{permanence test} - this must be established for a minimum prescribed
duration, twelve months with respect to a building site or construction or
installation project, [article 5(3)]\cite{oecd17_model_tax}, and
\item \textbf{business-activities test} lastly personally of the enterprise or agents on
their behalf, must conduct business activities at said place.
\end{itemize}

\subsection{Explain whether the following constitutes a PE on the Internet: a website; a server; a dependent agent. Refer to authoritative sources. (20)}
\label{sec:orgdeda852}

Adopting an antiquated view one may traditionally describe a \textbf{website} as a
collection of software and electronic data messages or constructs of electronic
data, as per [section 1]\cite{rsa02_elect_comm_trans_act}. From such a description
one `could' argue that a website is intangible and as such cannot alone
constitute a \emph{place of business} and thus failing to qualify as a (PE)
[paragraph 124]\cite{oecd17_model_tax}. Whereas on
the other hand, a \textbf{server} through which the website is stored and accessible is
a piece of machinery or equipment having a physical location and may thus
constitute a \emph{fixed place of business} as per the commentary to article 5
[paragraph 123]\cite{oecd17_model_tax}, and in turn a (PE).\\

Further distinction is required between the enterprises operating the server and
those that carry on business through the websites hosted on said servers.
Should an enterprise operating a website, also have at its own
disposal\footnote{Through ownership or lease agreement.}, use and operation of the
server hosting said website, then pending the other requirements the place where
that server is located could constitute a (PE) [paragraph
124]\cite{oecd17_model_tax}, provided it is fixed in that location for a
sufficient period of time [paragraph 125]\cite{oecd17_model_tax}.\\

In addition to the already mentioned reservations in establishing whether an
enterprise operating a website and/or server constitute a (PE), one must also
consider whether:
\begin{itemize}
\item the business of an enterprise is wholly or partially carried on through such
equipment needs to be examined on a case-by-case basis, [paragraph 126]\cite{oecd17_model_tax},
\item the computer equipment is automated and that no personnel of that enterprise
are required at the location [section 127]\cite{oecd17_model_tax},
\item no (PE) may exist where the computer equipment provides preparatory or
auxiliary services, for example the advertising goods and services, supplying
information, [paragraph 128]\cite{oecd17_model_tax}, unless these activities are
themselves essential and significant core functions of the enterprise [129]\cite{oecd17_model_tax},
\end{itemize}



During the course of 2015, OECD\footnote{Organization for Economic Co-operation and
Development} published the final version of their BEPS\footnote{G20 Base Erosion and
Profit Shifting} Project \cite{oecd15_action7}, with the aim of addressing tax
avoidance strategies that exploit the use of commissionaire arrangements to
avoid [article 5(5)]\cite{oecd17_model_tax}, and reliance of specific activity
exemptions [article 5(4)]\cite{oecd17_model_tax}. As per the provisions of
[article 5(5)]\cite{oecd17_model_tax} a person acting on behalf of a foreign
enterprise, then that enterprise shall be deemed to be a (PE) within a
Contracting State, despite not necessarily having a \emph{`fixed place-of-business'},
provided said person has and habitually exercises authority to conclude
contracts, or plays the principle role leading to the conclusion of contracts
that are routinely concluded without material modification by the enterprise, in
respect of activities undertaken for that enterprise [paragraph
31]\cite{oecd15_action7}. Such a person constitutes a \textbf{dependent agent} and in
such a capacity a dependent agent can indeed constitute a (PE),
\cite{oecd18_additional_guidance}.

\printbibliography
\end{document}